\section{Prism Refraction}

\label{sec:prism}
%% ============================================================================

Prism is QuasarSTM's primary contention-prevention layer. Block-STM~2.0
let conflicts happen and repaired them; Prism in 3.0 partitions the
ready frontier so that conflicting transactions execute in
\emph{different} slices and the STM rarely sees the same hot conflict
twice.

\subsection{Lane Conflict Matrix}

\begin{definition}[Lane conflict matrix]
\label{def:lane-conflict}
Given a ready frontier $F = \{T_1, \ldots, T_N\}$ with lane hint
vectors $\mathit{LH}(T_i)$, the lane conflict matrix
$M \in \{0,1\}^{N \times N}$ is defined by
\[
M_{ij} =
\begin{cases}
1 & \text{if } \exists L \in \mathit{LH}(T_i) \cap \mathit{LH}(T_j)
      \text{ s.t.\ at least one access is } \textsc{Write}, \\
0 & \text{otherwise.}
\end{cases}
\]
$M$ is symmetric; $M_{ii} = 0$.
\end{definition}

The matrix is computed in linear-arithmetic time on the GPU using
hashed lane bucket counts; full pairwise comparison is unnecessary.
The construction is deterministic.

\subsection{Slice Partition}

\begin{definition}[Refraction partition]
\label{def:refraction}
A refraction of $F$ is a partition
$F = S_1 \sqcup \cdots \sqcup S_q$ such that every slice $S_p$ falls
into one of three classes:
\begin{enumerate}[leftmargin=2em,nosep]
\item \textsc{DisjointWrite}: $\forall T_i, T_j \in S_p,\;
      M_{ij} = 0$.
\item \textsc{ReadOnly}: every $T_i \in S_p$ has only $\textsc{Read}$
      lane hints (and no writes will alter lane clocks).
\item \textsc{HotIsolated}: $S_p = \{T : L^* \in \mathit{LH}(T)\}$ for
      a single hot lane $L^*$; the slice is processed under a Tier-3
      reducer or a serialised lane mode.
\end{enumerate}
\end{definition}

\begin{algorithm}[H]
\caption{Compute refraction of frontier $F$}
\label{alg:refraction}
\begin{algorithmic}[1]
\State Bucket transactions by lane hint into a per-lane occupancy table.
\State Mark hot lanes $L^*$ where occupancy exceeds the hot-lane
       threshold $\theta_{\mathrm{hot}}$ or where
       $\mathit{LaneClock}[L^*].\mathit{hotness}$ exceeds
       $\theta_{\mathrm{hotness}}$.
\State \textbf{for each} hot lane $L^*$: emit \textsc{HotIsolated}
       slice $S_p \gets \{T : L^* \in \mathit{LH}(T)\}$ and remove
       these transactions from $F$.
\State Greedy-pack remaining transactions into \textsc{DisjointWrite}
       slices using a deterministic colouring of the (residual) lane
       conflict graph.
\State Emit a single \textsc{ReadOnly} slice for purely-read
       transactions.
\State \Return $\{S_1, \ldots, S_q\}$.
\end{algorithmic}
\end{algorithm}

\begin{lemma}[Refraction safety]
\label{lem:refraction-safe}
For any frontier $F$ and any refraction $\{S_1, \ldots, S_q\}$
satisfying Definition~\ref{def:refraction}, executing each slice
$S_p$ in parallel and serialising slices in the canonical order
respects $\preceq_{\mathrm{can}}$.
\end{lemma}

\begin{proof}[Sketch]
Within \textsc{DisjointWrite}, no two transactions touch a common
write lane, so no MVCC conflict is possible; any read-set under
unchanged lane clocks satisfies Tier-1.
Within \textsc{ReadOnly}, no writes alter visibility, and the
slice's transactions all observe the same pre-state.
Within \textsc{HotIsolated}, the registered Tier-3 reducer or
serialised-lane policy produces the canonical-order outcome.
Across slices, the canonical order is respected by sequential slice
dispatch.
\end{proof}

\subsection{Refraction Determinism}

\begin{invariant}[Refraction determinism]
\label{inv:refraction-deterministic}
For any frontier $F$ and any pair of executions of
Algorithm~\ref{alg:refraction} on the same $F$ and the same lane
hint vectors, the produced partition is bit-identical (up to slice
ordering, which is itself fixed by canonical-index lexicography).
\end{invariant}

This is mandatory: refraction sits on the round-result hash path.
A non-deterministic refraction would propagate to
\texttt{execution\_root} and break consensus.

\subsection{Why Refraction Is the Performance Primitive}

\paragraph{Same hot conflict twice.}
Without refraction, a frontier with 64 transactions all writing to
the AMM reserve produces $\binom{64}{2} = 2016$ pairwise conflicts on
the same lane. The validator detects them, the repair queue
absorbs them, and incarnation counters climb. With refraction, all
64 transactions land in a single \textsc{HotIsolated} slice and are
processed by a registered AMM reducer that yields exactly one
deterministic post-state.

\paragraph{Lane-fast hit rate.}
Refraction is the precondition for the Tier-1 hit-rate target
($> 95\%$ on disjoint-lane workloads, \S\ref{sec:performance}).
Without partitioning, any single hot transaction in a frontier
disturbs every co-resident transaction's lane snapshot.

\paragraph{Telemetry feedback loop.}
Hot-lane detection feeds the next round's refraction. A lane that
caused $> \theta_{\mathrm{hot}}$ conflicts in round $r$ is promoted
to \textsc{HotIsolated} for round $r+1$ \emph{before} the conflicts
recur. The control loop closes within one Quasar round.


%% ============================================================================
